\chapter{Apprentissage par renforcement~: fondements}\label{ch:ml:reinforcement-learning-foundations}

Un agent apprend à vendre un bloc d'actions par essais et erreurs dans un simulateur et bat le calendrier du desk de 0.9 point de base
par lot. Sur le marché dont le simulateur a été construit, il perd 0.6 point de base face au même calendrier. Le simulateur
comportait un détail faux~: une liquidité qui alternait à chaque période au lieu de persister, et l'agent avait appris à l'exploiter~:
il attendait chaque période de faible liquidité pour la période profonde que le simulateur livrait toujours ensuite. L'\omterm{def:ml:reinforcement-learning-foundations:mdp}{apprentissage par renforcement} est le contrôle optimal
appris à partir de l'expérience plutôt que résolu à partir d'un modèle (livre~4, chapitre~9)~; ce qu'il apprend, c'est
l'environnement dans lequel il a été entraîné. Ce chapitre travaille sur un problème assez petit pour être résolu exactement, afin que chaque apprenant puisse être
mesuré par rapport à la vérité~: les méthodes de valeur, les méthodes de \omterm{def:ml:reinforcement-learning-foundations:mdp}{politique}, les bandits, l'évaluation d'une nouvelle \omterm{def:ml:reinforcement-learning-foundations:mdp}{politique} à partir d'anciens journaux,
et l'écart entre un simulateur et un marché. Le chapitre~18 les applique à l'exécution et au market making.

\section{Processus de décision markoviens et équations de Bellman}

\begin{definition}[Apprentissage par renforcement, processus de décision markovien, politique, récompense, récompense cumulée]\label{def:ml:reinforcement-learning-foundations:mdp}
L'\emph{apprentissage par renforcement}\index{apprentissage par renforcement} apprend à agir à partir des conséquences des actions. Un
\emph{processus de décision markovien}\index{processus de décision markovien} (MDP) est un ensemble d'états $x$, d'actions $u$, de probabilités de transition
$\P(x_{t+1}\mid x_t, u_t)$ et une \emph{récompense}\index{récompense} $g_{t+1}$ reçue après chaque action~; la
\emph{récompense cumulée}\index{récompense cumulée} à partir de l'instant $t$ est $G_t = \sum_{k\ge0}\gamma^kg_{t+k+1}$, avec un facteur d'actualisation
$\gamma\in(0,1]$ ($\gamma = 1$ sur les horizons finis de ce chapitre). Une \emph{politique}\index{politique} $\pi(u\mid x)$ donne
la probabilité de chaque action dans chaque état.
\end{definition}

\begin{definition}[Fonction de valeur d'action, équation de Bellman]\label{def:ml:reinforcement-learning-foundations:bellman}
La \emph{fonction de valeur d'action}\index{fonction de valeur d'action} d'une \omterm{def:ml:reinforcement-learning-foundations:mdp}{politique} est $Q^\pi(x, u) = \E_\pi[G_t\mid x_t = x, u_t = u]$,
l'espérance de la \omterm{def:ml:reinforcement-learning-foundations:mdp}{récompense cumulée} lorsqu'on prend $u$ en $x$ puis qu'on suit $\pi$ ensuite~; la fonction valeur du livre~4 est $V^\pi(x)
= \sum_u\pi(u\mid x)Q^\pi(x, u)$. L'\emph{équation de Bellman}\index{equation de Bellman@équation de Bellman} est la récursion $Q^\pi(x, u) =
\E[g_{t+1} + \gamma V^\pi(x_{t+1})\mid x, u]$~; pour la \omterm{def:ml:reinforcement-learning-foundations:mdp}{politique} optimale, $Q^*(x, u) = \E[g_{t+1} + \gamma\max_{u'}Q^*(x_{t+1},
u')\mid x, u]$.
\end{definition}

Le problème du chapitre consiste à vendre 10 lots en 10 périodes. L'état est la période, le nombre de lots restants et la liquidité, profonde ou
faible, qui persiste d'une période à la suivante avec une probabilité de 0.8. Vendre $u$ lots coûte $\eta u^2$ points de base
d'impact temporaire, avec $\eta = 0.5$ lorsque la liquidité est profonde et 2 lorsqu'elle est faible~; garder $q$ lots après une vente coûte $0.2
q^2$ au titre du risque de mouvement du prix~; ce qui reste est vendu à la dernière période. La \omterm{def:ml:reinforcement-learning-foundations:mdp}{récompense} est l'opposé du coût. Avec 242
états, l'induction arrière sur l'\omterm{def:ml:reinforcement-learning-foundations:bellman}{équation de Bellman} (voir le \cref{lst:ml:rl:solve}) le résout exactement~: la \omterm{def:ml:reinforcement-learning-foundations:mdp}{politique} optimale coûte
2.59 points de base par lot (risque inclus), en vendant 5 lots d'un coup si la liquidité est profonde et 2 si elle est faible. Le calendrier
du desk, le meilleur qui ignore la liquidité, coûte 3.27~; vendre un lot par période coûte 6.76, à cause du risque de
détenir le bloc.

\section{Méthodes de valeur~: différences temporelles et Q-learning}

\begin{definition}[Apprentissage par différence temporelle, Q-learning]\label{def:ml:reinforcement-learning-foundations:td}
L'\emph{apprentissage par différence temporelle}\index{apprentissage par différence temporelle} met à jour une estimation de valeur vers une cible construite
à partir de la \omterm{def:ml:reinforcement-learning-foundations:mdp}{récompense} suivante et de l'estimation courante de la valeur de l'état suivant, $V(x_t)\leftarrow V(x_t) + \eta_\alpha\delta_t$ avec
$\delta_t = g_{t+1} + \gamma V(x_{t+1}) - V(x_t)$, sans attendre la fin de l'épisode. Le \emph{Q-learning}\index{Q-learning}
l'applique aux valeurs d'action avec la cible $g_{t+1} + \gamma\max_{u'}Q(x_{t+1}, u')$, ce qui apprend les valeurs de la \omterm{def:ml:reinforcement-learning-foundations:mdp}{politique} optimale
tout en agissant autrement (Watkins et Dayan, 1992)~; SARSA utilise l'action effectivement choisie ensuite, et apprend les valeurs de la
\omterm{def:ml:reinforcement-learning-foundations:mdp}{politique} qu'il suit.
\end{definition}

Le \Cref{lst:ml:rl:qlearning} est un \omterm{def:ml:reinforcement-learning-foundations:td}{Q-learning} tabulaire avec un pas d'apprentissage de 0.1 et une exploration $\varepsilon$-gloutonne (une action aléatoire
une fois sur dix). La \Cref{fig:ml:rl:curves} montre l'écart de valeur de la \omterm{def:ml:reinforcement-learning-foundations:mdp}{politique} apprise, mesuré exactement, en fonction du nombre
d'épisodes, moyenné sur trois graines. Après 1\,000 épisodes, la \omterm{def:ml:reinforcement-learning-foundations:mdp}{politique} gloutonne du \omterm{def:ml:reinforcement-learning-foundations:td}{Q-learning} est encore à 1.41 points de base par lot de
l'optimum~; après 10\,000, elle n'en est plus qu'à 0.044. SARSA, qui apprend la valeur de sa propre \omterm{def:ml:reinforcement-learning-foundations:mdp}{politique} d'exploration, est plus prudent au début
(0.28 à 1\,000 épisodes) et termine à 0.051.

\section{Méthodes de politique~: gradients et acteur--critique}

\begin{definition}[Gradient de politique, méthode acteur--critique]\label{def:ml:reinforcement-learning-foundations:pg}
Une méthode de \emph{gradient de politique}\index{gradient de politique} ajuste les paramètres $\theta$ d'une \omterm{def:ml:reinforcement-learning-foundations:mdp}{politique} stochastique dans la direction
$\E[G_t\nabla_\theta\log\pi_\theta(u_t\mid x_t)]$, le gradient de la \omterm{def:ml:reinforcement-learning-foundations:mdp}{récompense cumulée} espérée (REINFORCE~; Williams, 1992)~; une
référence soustraite à $G_t$ réduit la variance sans la biaiser. Une
\emph{méthode acteur--critique}\index{méthode acteur--critique} remplace $G_t$ par l'erreur de différence temporelle d'une fonction de valeur apprise (le critique), de sorte que la \omterm{def:ml:reinforcement-learning-foundations:mdp}{politique}
(l'acteur) apprend à chaque pas.
\end{definition}

REINFORCE avec une \omterm{def:ml:reinforcement-learning-foundations:mdp}{politique} softmax et une référence par état apprend lentement et de façon bruitée~: 0.31 point de base de l'optimum après
1\,000 épisodes, 0.067 après 10\,000, et avec un pas d'apprentissage plus grand, il s'est fixé sur une \omterm{def:ml:reinforcement-learning-foundations:mdp}{politique} moins bonne qu'il n'a jamais quittée. La
méthode acteur-critique du \cref{lst:ml:rl:ac} apprend le plus vite ici (0.11 à 1\,000 épisodes, 0.009 à 10\,000). Les méthodes de \omterm{def:ml:reinforcement-learning-foundations:mdp}{politique} prennent de l'importance
lorsque les actions sont continues ou nombreuses, situation où un maximum sur les actions est coûteux, et lorsqu'une \omterm{def:ml:reinforcement-learning-foundations:mdp}{politique} stochastique est elle-même le
but~; les versions profondes remplacent les tableaux par des réseaux (Mnih et coauteurs, 2015, pour les méthodes de valeur).

\begin{omfigure}
\begin{tikzpicture}
\begin{loglogaxis}[width=11cm,height=6cm,xlabel={épisodes d'entraînement},ylabel={écart à l'optimum (pb par lot)},
  tick label style={font=\small},label style={font=\small},xmin=80,xmax=12500,ymin=0.005,ymax=3,grid=major,
  grid style={gray!25},legend columns=4,legend style={font=\footnotesize,at={(0.5,-0.25)},anchor=north,draw=none,
  fill=none,/tikz/every even column/.append style={column sep=6pt}},table/col sep=comma]
\addplot[omThm,very thick,mark=*] table[x=episodes,y=q]{figdata/ml/17-reinforcement-learning-foundations/curves.csv};
\addplot[omProp,very thick,mark=square*] table[x=episodes,y=sarsa]{figdata/ml/17-reinforcement-learning-foundations/curves.csv};
\addplot[gray,very thick,mark=triangle*] table[x=episodes,y=reinforce]{figdata/ml/17-reinforcement-learning-foundations/curves.csv};
\addplot[omDef,very thick,mark=diamond*] table[x=episodes,y=ac]{figdata/ml/17-reinforcement-learning-foundations/curves.csv};
\legend{Q-learning,SARSA,REINFORCE,acteur-critique}
\end{loglogaxis}
\end{tikzpicture}

\omcaption{Écart de valeur à l'optimum de la programmation dynamique pour la \omterm{def:ml:reinforcement-learning-foundations:mdp}{politique} apprise après un certain nombre d'épisodes (la \omterm{def:ml:reinforcement-learning-foundations:mdp}{politique} gloutonne
pour le \omterm{def:ml:reinforcement-learning-foundations:td}{Q-learning} et SARSA, la \omterm{def:ml:reinforcement-learning-foundations:mdp}{politique} softmax pour les autres), moyenne sur trois graines. Données~:
\texttt{ml\_rl.learning\_curves}.}\label{fig:ml:rl:curves}
\end{omfigure}

\section{Bandits et exploration}

\begin{definition}[Arbitrage exploration--exploitation, bandit manchot à plusieurs bras, bandit contextuel]\label{def:ml:reinforcement-learning-foundations:bandit}
L'\emph{arbitrage exploration--exploitation}\index{arbitrage exploration--exploitation} est le choix entre l'action qui
semble la meilleure maintenant et une action qui en apprend davantage. Un \emph{bandit manchot à plusieurs bras}\index{bandit manchot à plusieurs bras} (multi-armed bandit) est le MDP à un
seul état~: à chaque tour un bras est choisi et sa \omterm{def:ml:reinforcement-learning-foundations:mdp}{récompense} tirée~; la perte par rapport au fait de toujours jouer le meilleur bras est le regret. Un
\emph{bandit contextuel}\index{bandit contextuel} observe un contexte (l'ordre, le marché) avant de choisir, et apprend un
modèle de \omterm{def:ml:reinforcement-learning-foundations:mdp}{récompense} par bras~; ses actions ne modifient pas les contextes futurs.
\end{definition}

Cinq algorithmes d'exécution font économiser en moyenne 0, 0.5, 1.0, 1.2 et 1.5 points de base par ordre, avec 5 points de base de bruit par
ordre. Sur 10\,000 ordres (voir la \cref{fig:ml:rl:regret}), l'$\varepsilon$-glouton avec $\varepsilon = 0.1$ perd 1\,727 points de base
d'économies cumulées par rapport au fait de toujours utiliser le meilleur algorithme, et l'utilise dans 78~\% des mille derniers ordres~; avec
$\varepsilon = 0.01$, il explore trop peu et se verrouille sur le mauvais algorithme dans certains essais (2\,026, 65~\%). La règle de la
borne de confiance supérieure, qui choisit le bras de moyenne la plus élevée augmentée de $c$ écarts-types du bruit multipliés par $\sqrt{\ln t/n_a}$
(Auer, Cesa-Bianchi et Fischer, 2002), perd 1\,325 avec $c = \sqrt2$ et 944 avec $c = 1$ (83~\% et 87~\%)~: elle explore les
bras dont elle doute, et non au hasard. L'essentiel de la perte est inévitable~: un bras inférieur de 0.3 point de base au meilleur, sous
5 points de base de bruit, demande environ deux mille ordres sur chacun pour être distingué à deux erreurs types.

\begin{omfigure}
\begin{tikzpicture}
\begin{axis}[width=11cm,height=6cm,xlabel={ordres},ylabel={regret cumulé (pb)},tick label style={font=\small},
  label style={font=\small},xmin=0,xmax=10000,ymin=0,grid=major,grid style={gray!25},
  xtick={0,2000,4000,6000,8000,10000},xticklabel style={/pgf/number format/1000 sep={\,}},scaled x ticks=false,legend cell align=left,
  legend style={font=\footnotesize,at={(0.02,0.98)},anchor=north west,fill=white,draw=none},table/col sep=comma]
\addplot[omThm,very thick] table[x=t,y=eps1]{figdata/ml/17-reinforcement-learning-foundations/regret.csv};
\addplot[omThm!50,very thick,dashed] table[x=t,y=eps01]{figdata/ml/17-reinforcement-learning-foundations/regret.csv};
\addplot[omDef!60,very thick,dashed] table[x=t,y=ucb141]{figdata/ml/17-reinforcement-learning-foundations/regret.csv};
\addplot[omDef,very thick] table[x=t,y=ucb1]{figdata/ml/17-reinforcement-learning-foundations/regret.csv};
\legend{$\varepsilon$-glouton 0.1,$\varepsilon$-glouton 0.01,UCB $c = \sqrt2$,UCB $c = 1$}
\end{axis}
\end{tikzpicture}

\omcaption{Regret cumulé de quatre règles d'exploration choisissant parmi cinq algorithmes d'exécution, moyenne de 20 essais. Données~:
\texttt{ml\_rl.bandits}.}\label{fig:ml:rl:regret}
\end{omfigure}

\section{Évaluation hors politique et le problème du simulateur}

\begin{definition}[Évaluation hors politique, estimateur doublement robuste]\label{def:ml:reinforcement-learning-foundations:ope}
L'\emph{évaluation hors politique}\index{evaluation hors politique@évaluation hors politique} (off-policy evaluation) estime la valeur d'une \omterm{def:ml:reinforcement-learning-foundations:mdp}{politique} cible à partir d'épisodes générés par
une autre, la \omterm{def:ml:reinforcement-learning-foundations:mdp}{politique} comportementale, typiquement en repondérant chaque épisode par le rapport entre les probabilités des deux \omterm{def:ml:reinforcement-learning-foundations:mdp}{politiques} pour les
actions effectivement prises (échantillonnage préférentiel, livre~4, chapitre~26). L'\emph{estimateur doublement robuste}\index{estimateur doublement robuste} (doubly robust estimator)
ajoute à l'estimation du modèle de la valeur de la cible les erreurs du modèle sur les \omterm{def:ml:reinforcement-learning-foundations:mdp}{récompenses} journalisées, pondérées par les rapports de probabilités~; elle est
sans biais si les probabilités ou le modèle sont justes, et a une variance plus faible que l'échantillonnage préférentiel lorsque le modèle
est proche de la réalité (Dudík, Langford et Li, 2011).
\end{definition}

Un desk ne peut guère essayer une nouvelle \omterm{def:ml:reinforcement-learning-foundations:mdp}{politique} sur de vrais ordres avant d'en connaître la valeur. Ici, les journaux du desk proviennent de son
calendrier, avec des écarts aléatoires dans 30~\% des cas (coût 3.80)~; la cible est la \omterm{def:ml:reinforcement-learning-foundations:mdp}{politique} optimale avec des écarts dans 10~\% des
cas, de sorte que chaque action qu'elle peut prendre a parfois été journalisée (coût réel 2.77). Le \Cref{tab:ml:rl:ope} compare les estimateurs
sur 100 jeux de 300 épisodes journalisés~; le modèle fourni à l'\omterm{def:ml:reinforcement-learning-foundations:ope}{estimateur doublement robuste} surestime de moitié
l'impact en liquidité faible et rend la liquidité trop variable.

\begin{table}[htbp]
\centering\small
\begin{tabular}{@{}lccc@{}}
\toprule
estimateur & biais & écart-type & erreur quadratique moyenne\\
\midrule
échantillonnage préférentiel par décision & 0.20 & 1.75 & 1.76\\
échantillonnage préférentiel pondéré & 0.49 & 0.63 & 0.80\\
modèle seul & 0.30 & 0 & 0.30\\
doublement robuste & $-0.03$ & 0.29 & 0.29\\
\bottomrule
\end{tabular}
\caption{Estimations du coût de la \omterm{def:ml:reinforcement-learning-foundations:mdp}{politique} cible (valeur réelle 2.77 pb par lot) à partir de 300 épisodes journalisés de la
\omterm{def:ml:reinforcement-learning-foundations:mdp}{politique} comportementale du desk, sur 100 répétitions (pb par lot). Données~: \texttt{ml\_rl.off\_policy}.}\label{tab:ml:rl:ope}
\end{table}

L'échantillonnage préférentiel est sans biais mais inutilisable seul~: des produits de dix rapports de probabilités font dominer quelques épisodes. Pondérer
les épisodes pour que leurs poids somment à un revient à échanger de la variance contre du biais. Le modèle seul est précis mais faux~; l'\omterm{def:ml:reinforcement-learning-foundations:ope}{estimateur doublement robuste} conserve
la précision du modèle et élimine l'essentiel de son biais (\cref{lst:ml:rl:dr}).

\begin{definition}[Écart simulation--réalité]\label{def:ml:reinforcement-learning-foundations:simreal}
L'\emph{écart simulation--réalité}\index{ecart simulation--realite@écart simulation--réalité} (sim-to-real gap) est la différence entre la performance d'une \omterm{def:ml:reinforcement-learning-foundations:mdp}{politique} dans le simulateur
où elle a été entraînée et dans l'environnement réel. Il est le plus grand pour les \omterm{def:ml:reinforcement-learning-foundations:mdp}{politiques} qui exploitent ce que le simulateur modélise mal,
c'est-à-dire ce que fait un optimiseur avec toute erreur qui rapporte.
\end{definition}

L'agent de l'accroche est un \omterm{def:ml:reinforcement-learning-foundations:td}{Q-learning} entraîné sur 10\,000 épisodes dans un simulateur identique au marché, sauf que la liquidité
alterne à chaque période, comme elle le faisait selon un schéma calendaire dans les données sur lesquelles le simulateur a été calibré (livre~7, chapitre~19). Dans
le simulateur, il coûte 2.77 points de base par lot contre 3.66 pour le desk. Sur le marché, il coûte 3.83 contre
3.27 pour le desk et 2.59 pour l'optimum~: en attendant une liquidité profonde qui ne vient pas, il garde son risque plus longtemps et vend dans des marchés minces
en fin de parcours. Le calendrier du desk, qui ne supposait rien sur la liquidité, n'a rien perdu lorsque l'hypothèse s'est révélée fausse.

\begin{method}[Avant de faire confiance à une politique apprise]\label{met:ml:rl:method}
\begin{enumerate}
\item Résoudre exactement la plus petite version du problème et vérifier que l'apprenant la retrouve.
\item Estimer la valeur de la nouvelle \omterm{def:ml:reinforcement-learning-foundations:mdp}{politique} à partir de journaux réels par \omterm{def:ml:reinforcement-learning-foundations:ope}{évaluation hors politique}, avec un \omterm{def:ml:reinforcement-learning-foundations:ope}{estimateur doublement robuste} et les
probabilités journalisées de la \omterm{def:ml:reinforcement-learning-foundations:mdp}{politique} comportementale~; les consigner pour chaque décision.
\item Évaluer la \omterm{def:ml:reinforcement-learning-foundations:mdp}{politique} dans des simulateurs qui diffèrent de celui d'entraînement par les hypothèses qu'elle pourrait exploiter (persistance,
impact, exécutions), et privilégier les \omterm{def:ml:reinforcement-learning-foundations:mdp}{politiques} dont la valeur est stable entre eux.
\item Comparer avec le calendrier le plus simple, qui suppose le moins, et déployer avec des limites (chapitre~18).
\end{enumerate}
\end{method}

\section{Tutoriel~: apprendre à vendre}\label{tut:ml:reinforcement-learning-foundations:tut}

\begin{tutorial}
\textbf{Objectif.} Résoudre exactement le MDP de liquidation, l'apprendre de quatre façons, faire tourner les bandits, évaluer une \omterm{def:ml:reinforcement-learning-foundations:mdp}{politique} à partir de journaux et entraîner
dans un simulateur erroné. \textbf{État final~:} \cref{fig:ml:rl:curves,fig:ml:rl:regret}, \cref{tab:ml:rl:ope}.
\begin{tutsteps}
\item \textbf{Induction arrière.}
\omcode{code/firm/rlcore/firm_rlcore.py}{36}{46}{La récursion d'optimalité de Bellman sur un horizon fini.}\label{lst:ml:rl:solve}
\item \textbf{\omterm{def:ml:reinforcement-learning-foundations:td}{Q-learning} et SARSA.}
\omcode{code/firm/rlcore/firm_rlcore.py}{128}{148}{Q-learning tabulaire (et SARSA) avec une exploration epsilon-gloutonne.}\label{lst:ml:rl:qlearning}
\item \textbf{Acteur-critique.}
\omcode{code/firm/rlcore/firm_rlcore.py}{210}{229}{Un acteur-critique à un pas avec une politique softmax.}\label{lst:ml:rl:ac}
\item \textbf{Évaluation doublement robuste.}
\omcode{code/firm/rlcore/firm_rlcore.py}{290}{301}{L'estimateur doublement robuste par décision.}\label{lst:ml:rl:dr}
\item \textbf{Exécuter} \texttt{ml\_rl.benchmarks()}, \texttt{learning\_curves()}, \texttt{bandits()}, \texttt{off\_policy()},
\texttt{sim\_to\_real()} et \texttt{fig\_rl.py}.
\end{tutsteps}
\textbf{Pour aller plus loin.} Faire en sorte que les ventes de l'agent épuisent la liquidité (le paramètre \texttt{drain}) et que l'agent soit réentraîné~; donner à
l'\omterm{def:ml:reinforcement-learning-foundations:ope}{estimateur doublement robuste} le vrai modèle, puis un modèle moins bon.
\end{tutorial}

\section{Construction~: noyau d'apprentissage par renforcement}\label{bld:ml:reinforcement-learning-foundations:rlcore}

\begin{build}
\bfield{Rôle} \omterm{def:ml:reinforcement-learning-foundations:mdp}{Apprentissage par renforcement} testé contre des réponses exactes, et \omterm{def:ml:reinforcement-learning-foundations:mdp}{politiques} évaluées avant de trader.
\bfield{Interface} \texttt{FiniteMDP} avec \texttt{solve}, \texttt{evaluate}, \texttt{q\_values}~;
\texttt{liquidation\_mdp(T, Q, eta, p\_stay, phi, drain)}~; \texttt{MDPEnv} avec \texttt{reset(seed)}, \texttt{step(u)}~;
\texttt{q\_learning}, \texttt{sarsa}, \texttt{greedy}, \texttt{reinforce}, \texttt{actor\_critic}, \texttt{softmax\_policy}~;
\texttt{bandit(means, sd, n, rule, seed, eps, c)}~; \texttt{rollouts}, \texttt{ope\_is}, \texttt{ope\_wis}, \texttt{ope\_dr}.
\bfield{Règles} Tout apprenant est initialisé avec une graine et reproductible~; les \omterm{def:ml:reinforcement-learning-foundations:mdp}{politiques} apprises sont évaluées exactement là où le MDP est connu~;
les données journalisées portent les probabilités de la \omterm{def:ml:reinforcement-learning-foundations:mdp}{politique} comportementale.
\bfield{Tests d'acceptation} \texttt{code/firm/rlcore/tests/}~: l'induction arrière coïncide avec l'énumération exhaustive sur un MDP minuscule~;
l'évaluation exacte coïncide avec Monte Carlo~; le \omterm{def:ml:reinforcement-learning-foundations:td}{Q-learning} atteint l'optimum sur un petit problème~; l'échantillonnage préférentiel est sans biais
et l'\omterm{def:ml:reinforcement-learning-foundations:ope}{estimateur doublement robuste} a une dispersion plus faible, et il est exact avec un modèle parfait et des transitions déterministes~; le regret de l'UCB croît plus lentement que celui d'une \omterm{def:ml:reinforcement-learning-foundations:mdp}{politique} aléatoire.
\bfield{Pour aller plus loin} Approximation de fonction (un réseau pour $Q$, chapitre~18)~; LinUCB pour les \omterm{def:ml:reinforcement-learning-foundations:bandit}{bandits contextuels}~; \omterm{def:ml:reinforcement-learning-foundations:ope}{évaluation hors politique}
avec des probabilités comportementales estimées.
\end{build}

\begin{omsources}
\item R.~S.~Sutton et A.~G.~Barto, \emph{Reinforcement Learning: An Introduction}, 2\up{e} éd., MIT Press, 2018.
\item C.~J.~C.~H.~Watkins et P.~Dayan, ``Q-learning'', \emph{Machine Learning} 8, 1992.
\item R.~J.~Williams, ``Simple statistical gradient-following algorithms for connectionist reinforcement learning'',
\emph{Machine Learning} 8, 1992.
\item P.~Auer, N.~Cesa-Bianchi et P.~Fischer, ``Finite-time analysis of the multiarmed bandit problem'', \emph{Machine
Learning} 47, 2002.
\item M.~Dudík, J.~Langford et L.~Li, ``Doubly robust policy evaluation and learning'', arXiv:1103.4601, 2011.
\item N.~Jiang et L.~Li, ``Doubly robust off-policy value evaluation for reinforcement learning'', arXiv:1511.03722, 2016.
\item V.~Mnih et coauteurs, ``Human-level control through deep reinforcement learning'', \emph{Nature} 518, 2015.
\end{omsources}

\section{Exercices}

\begin{exercise}[$\star$]\label{exo:ml:reinforcement-learning-foundations:1}
Dans la dernière période, la totalité du reliquat doit être vendue. Écrire la valeur optimale de l'état (période 9, 3 lots,
liquidité faible) et celle de (période 9, 3 lots, liquidité profonde).
\end{exercise}

\begin{exercise}[$\star$]\label{exo:ml:reinforcement-learning-foundations:2}
Une mise à jour en \omterm{def:ml:reinforcement-learning-foundations:td}{Q-learning} a $Q(x, u) = -10$, une \omterm{def:ml:reinforcement-learning-foundations:mdp}{récompense} $-3$, $\max_{u'}Q(x', u') = -5$ et un pas d'apprentissage 0.1. Quelle est la nouvelle valeur de $Q(x,
u)$~? Quelle valeur SARSA utiliserait-il si l'action suivante effectivement prise avait $Q(x', u') = -8$~?
\end{exercise}

\begin{exercise}[$\star$]\label{exo:ml:reinforcement-learning-foundations:3}
Pourquoi le coût de la vente d'un lot par période (6.76) dépasse-t-il à ce point le calendrier du desk (3.27) dans ce problème~?
\end{exercise}

\begin{exercise}[$\star\star$]\label{exo:ml:reinforcement-learning-foundations:4}
Pourquoi SARSA devance-t-il le \omterm{def:ml:reinforcement-learning-foundations:td}{Q-learning} après 1\,000 épisodes et lui est-il inférieur après 10\,000~?
\end{exercise}

\begin{exercise}[$\star\star$]\label{exo:ml:reinforcement-learning-foundations:5}
Pourquoi l'échantillonnage préférentiel ordinaire présente-t-il ici un écart type aussi grand~? Estimer la taille typique du poids
d'un épisode.
\end{exercise}

\begin{exercise}[$\star\star$]\label{exo:ml:reinforcement-learning-foundations:6}
\emph{Trouver la faille.} «~Notre agent d'exécution en \omterm{def:ml:reinforcement-learning-foundations:mdp}{apprentissage par renforcement} a battu le TWAP de 2 points de base sur un an de simulation, avec une
$t$-statistique de 40, donc nous le déployons.~»
\end{exercise}

\begin{exercise}[$\star\star\star$]\label{exo:ml:reinforcement-learning-foundations:7}
\emph{Programmation.} Fixer \texttt{drain = 0.1} dans le marché réel (chaque lot vendu accroît de
0.1 la probabilité que la période suivante soit peu liquide) et comparer le coût optimal et la première action avec ceux du marché du chapitre. Expliquer la différence.
\end{exercise}

\begin{exercise}[$\star\star\star$]\label{exo:ml:reinforcement-learning-foundations:8}
Montrer que l'\omterm{def:ml:reinforcement-learning-foundations:ope}{estimateur doublement robuste} est sans biais lorsque le modèle est exact ($\hat Q = Q^{\pi_e}$), quels que soient les
poids d'importance, pour un problème à une étape.
\end{exercise}

\section{Problème~: Apprendre à vendre}

\begin{problem}[{Problème du week-end --- une politique ne vaut que par le monde qui la sous-tend}]\label{pb:ml:reinforcement-learning-foundations:1}
Le problème de liquidation du chapitre, les apprenants, les bandits, les journaux et le simulateur.

\medskip\noindent\textbf{Partie I --- Le problème exact.}
\begin{enumerate}
\item Décrire les états, les actions, les \omterm{def:ml:reinforcement-learning-foundations:mdp}{récompenses} et les transitions.
\item Quel est le coût de la \omterm{def:ml:reinforcement-learning-foundations:mdp}{politique} optimale, et que fait-elle en premier lieu en liquidité profonde et en liquidité faible~?
\item Combien coûtent le calendrier du desk et la vente d'un lot par période~?
\item Pourquoi l'induction arrière est-elle possible ici et ne l'est-elle pas sur un marché réel~?
\end{enumerate}

\medskip\noindent\textbf{Partie II --- L'apprendre.}
\begin{enumerate}[resume]
\item Quel écart à l'optimum laisse chaque apprenant après 1\,000 et après 10\,000 épisodes~?
\item Pourquoi REINFORCE peut-il se fixer sur une \omterm{def:ml:reinforcement-learning-foundations:mdp}{politique} moins bonne avec un pas plus grand~?
\item Qu'apporte la \omterm{def:ml:reinforcement-learning-foundations:pg}{méthode acteur--critique} par rapport à REINFORCE~?
\item À combien de jours de négociation correspondent 10\,000 épisodes, si un épisode est un ordre parent~?
\end{enumerate}

\medskip\noindent\textbf{Partie III --- Bandits et journaux.}
\begin{enumerate}[resume]
\item Quel regret accumule chaque règle d'exploration, et laquelle est la meilleure~?
\item Pourquoi $\varepsilon = 0.01$ fait-il moins bien que 0.1~?
\item Comparer les quatre estimations hors \omterm{def:ml:reinforcement-learning-foundations:mdp}{politique}.
\item Que doit journaliser le desk pour que l'\omterm{def:ml:reinforcement-learning-foundations:ope}{évaluation hors politique} soit possible~?
\end{enumerate}

\medskip\noindent\textbf{Partie IV --- Le verdict.}
\begin{enumerate}[resume]
\item Énoncer le \emph{résultat nommé}~: l'écart de valeur du \omterm{def:ml:reinforcement-learning-foundations:td}{Q-learning} à l'optimum de programmation dynamique après 10\,000 épisodes, ainsi que le
biais et la dispersion des trois estimateurs hors \omterm{def:ml:reinforcement-learning-foundations:mdp}{politique}.
\item Qu'a appris l'agent entraîné dans le mauvais simulateur, et combien cela a-t-il coûté~?
\item Comment l'aurait-on détecté avant le déploiement~?
\item Quand préférer un bandit à l'\omterm{def:ml:reinforcement-learning-foundations:mdp}{apprentissage par renforcement} complet pour un choix d'exécution~?
\item Pourquoi le calendrier du desk est-il robuste~?
\item Qu'est-ce qui change avec des prix et des tailles continus~?
\item Quel rôle doit jouer un simulateur dans la mise au point d'une \omterm{def:ml:reinforcement-learning-foundations:mdp}{politique} de trading~?
\item En une phrase~: qu'apprend un agent d'\omterm{def:ml:reinforcement-learning-foundations:mdp}{apprentissage par renforcement}~?
\end{enumerate}
\end{problem}

\section{Questions d'entretien}

\begin{interviewq}[$\star$ \iqroles{chercheur, ingénieur ML}]\label{iq:ml:reinforcement-learning-foundations:1}
Écrivez l'équation d'optimalité de Bellman pour la \omterm{def:ml:reinforcement-learning-foundations:bellman}{fonction de valeur d'action} et expliquez chacun de ses termes.
\end{interviewq}

\begin{interviewq}[$\star\star$ \iqroles{ingénieur ML}]\label{iq:ml:reinforcement-learning-foundations:2}
Quelle est la différence entre le \omterm{def:ml:reinforcement-learning-foundations:td}{Q-learning} et SARSA~? Quand cela importe-t-il~?
\end{interviewq}

\begin{interviewq}[$\star\star$ \iqroles{chercheur}]\label{iq:ml:reinforcement-learning-foundations:3}
Comment estimeriez-vous la valeur d'une nouvelle \omterm{def:ml:reinforcement-learning-foundations:mdp}{politique} d'exécution à partir des exécutions historiques, sans la mettre en œuvre~?
\end{interviewq}

\begin{interviewq}[$\star\star$ \iqroles{chercheur, trader}]\label{iq:ml:reinforcement-learning-foundations:4}
Vous devez choisir parmi cinq algorithmes de courtiers pour chaque ordre. Concevez la règle de choix.
\end{interviewq}

\begin{interviewq}[$\star\star$ \iqroles{ingénieur ML}]\label{iq:ml:reinforcement-learning-foundations:5}
Expliquez le théorème du \omterm{def:ml:reinforcement-learning-foundations:pg}{gradient de politique} et pourquoi une référence (baseline) aide.
\end{interviewq}

\begin{interviewq}[$\star\star\star$ \iqroles{chercheur}]\label{iq:ml:reinforcement-learning-foundations:6}
Un agent entraîné dans votre simulateur de marché bat tous les indices de référence. Énumérez ce qui pourrait le faire échouer en production réelle, et comment vous testeriez
chacun de ces points.
\end{interviewq}
